\section*{Introduction}

Symmetry is the organising principle of machine learning for atomistic systems. Networks whose internal features transform equivariantly under the symmetries of three-dimensional space have become the default architecture for predicting the properties of molecules and crystals, from interatomic potentials\cite{batzner2022nequip,musaelian2023allegro,batatia2022mace} to equivariant transformers\cite{liao2023equiformer,liao2024equiformerv2}, and the community describes them almost interchangeably as ``E(3)-equivariant''. The label conceals a fork. The Euclidean group E(3) contains improper operations: reflections and the spatial inversion $I\colon\, \mathbf{r} \mapsto -\mathbf{r}$. A network can be exactly equivariant to all rotations, SO(3), while carrying no information about parity at all --- and a substantial class of deployed models, including the eSCN lineage\cite{passaro2023escn} and EquiformerV2, is of exactly this kind. The distinction is invisible at the configuration surface: such models expose a maximum angular degree, not a parity flag, and the newest generation of the lineage continues to advertise itself by its rotational symmetry group alone\cite{equiformerv3}. Whether a model respects improper operations is, in practice, an implementation detail that neither model cards nor benchmarks report.

How much symmetry to build into these models is an active debate, but it has been argued almost entirely on the rotation axis. There, the stakes are accuracy: models that are only approximately rotation-equivariant have been found to suffer negligible consequences in realistic bulk simulation\cite{langer2024broken}, and unconstrained architectures can learn rotational symmetry from data. This paper concerns the other axis, where the stakes are different in kind. A model that is slightly wrong about a rotation makes a slightly wrong prediction; a model that is blind to parity can make predictions that are not approximately wrong but physically impossible.

The physics supplies an unusually clean instrument for making this precise. Neumann's principle\cite{nye1985crystals} requires every macroscopic property tensor of a crystal to be invariant under each symmetry operation of the crystal's point group. The piezoelectric tensor $e_{ijk}$, which couples polarisation to strain, is a parity-odd rank-3 tensor: under inversion it changes sign. For a centrosymmetric crystal --- one whose space group contains the inversion --- invariance and sign reversal are simultaneously required, so the tensor must equal its own negative. The consequence is an algebraic identity, not an empirical trend: the piezoelectric tensor of every centrosymmetric crystal is exactly zero. Since 92 of the 230 space groups are centrosymmetric, this yields a large, label-free test population on which any nonzero prediction is known, with certainty and without reference data, to be impossible.

Whether an equivariant network can violate this identity is decided by its feature algebra. An O(3)-equivariant network attaches a parity label to every internal feature, and the composition rules of those labels force any parity-odd output to vanish identically on a centrosymmetric input: the zero is structural, holding at any parameter values, including random initialisation. An SO(3)-equivariant network is the same construction with the parity labels erased. It remains exactly rotation-equivariant and loses little accuracy on standard benchmarks, but nothing in it forbids a symmetry-forbidden output. The question this paper asks is what that missing guarantee costs on real crystals, in trained, deployed-scale models.

Two adjacent lines of work make the question timely without answering it. A literature on symmetry breaking in equivariant networks starts from the observation, rooted in Curie's principle\cite{curie1894symmetry}, that an equivariant model cannot produce outputs less symmetric than its inputs, and develops relaxations that permit such outputs when the physics demands them\cite{smidt2021symmetry,xie2024symmetrybreaking,kaba2023symmetry,prob2025symmetry}. We probe the complement: for a parity-odd tensor on a centrosymmetric crystal there is no symmetry breaking to model, the answer is exactly zero, and a model that escapes the constraint is not flexible but wrong. Separately, a growing family of crystal-tensor predictors enforces the correct symmetry at the output, by space-group-informed enforcement modules, canonicalisation, or symmetry-adapted decompositions\cite{yan2024gmtnet,goectp2024,lu2025eatgnn,piezotensornet2024,ceitnet2026}. These are readout-level repairs; whether the features beneath them can represent the constraint is a separate and prior question, and it becomes urgent as property heads are increasingly attached to embeddings of pretrained universal potentials\cite{hung2024dielectric}, where the backbone's symmetry group is fixed long before any downstream target is chosen.

In this work we measure what parity labels are worth. We build matched pairs of three architectures --- NequIP, Allegro, and MACE --- that are identical in every respect except the parity labelling of their features, add EquiformerV2 unmodified as a deployed rotation-only representative, and evaluate all of them on a large spglib-verified population of centrosymmetric insulators where the true piezoelectric tensor vanishes identically. We find that the parity-labelled arms predict zero for every crystal while the rotation-only arms predict tensors as large as genuine piezoelectric responses for roughly nine crystals in ten; that this failure cannot be trained away by zero-labelled data, loss reweighting, or test-time symmetrisation; that the crystals the rotation-only models do get right are precisely those where rotational symmetry alone forbids the tensor; and that the parity-labelled guarantee is differentiable, tracking the physics of symmetry breaking. An audit of fifteen released architectures shows that most cannot structurally forbid impossible parity-odd predictions. The symmetry group of a model's features, not its training data, decides whether physically impossible predictions are possible.
